\documentclass[../../main]{subfiles}

\ifSubfilesClassLoaded{%
    \zexternaldocument*{1.Categories-morphisms-and-duality}%
    \zexternaldocument*{2.Functors-and-natural-transformations}%
    \zexternaldocument*{3.Limits-and-colimits}
}{}


\begin{document}
\chapter{More on functors}

\section{Representable functors}
\begin{definition}\label{def:representable-functor}
  A \hyperref[def:functor]{functor} \(F: \category{C} \to \Set\) is \textbf{representable} if there exists an object \(X \in \category{C}\) such that
  \begin{equation*}
    F \cong \Hom_{\category{C}}(X, -)
  \end{equation*}
  where \(\Hom_{\category{C}}(X, -)\) is the \hyperref[ex:hom-functor]{covariant Hom functor}.
\end{definition}

\begin{definition}[Concrete category]\label{def:concrete-category}
  A \hyperref[def:category]{category} \(\category{C}\) is \textbf{concrete} if there exists a \hyperref[def:functor-types]{faithful functor} \(U: \category{C} \to \Set\). Such a functor is often called a \textbf{forgetful functor}.
\end{definition}

\begin{lemma}[{Yoneda - \cite[Theorem 2.2.4]{riehlCategoryTheoryContext}}]\label{lemma:yoneda}
  Let \(\category{C}\) be a \hyperref[def:category]{locally small category}, \(X \in \category{C}\) an object, and \(F: \category{C} \to \Set\) a \hyperref[def:functor]{functor}. There is a bijection:
  \begin{equation*}
    \Phi: \Hom(\Hom(X, -), F) \isoarrow F(X)
  \end{equation*}
  given by \(\alpha \mapsto \alpha_{X}(\id_{X})\). This bijection is natural in both \(X\) and \(F\).
\end{lemma}

\section{Adjunctions}

\begin{definition}[Adjoint functors]\label{def:adjoint-functors}
  Let \(\category{C}\), \(\category{D}\) be two \hyperref[def:category]{locally small categories} and let \(L: \category{C} \to \category{D}\) and \(R: \category{D} \to \category{C}\) be two \hyperref[def:functor]{functors}. We say that \(L\) is \textbf{left adjoint} to \(R\), and \(R\) is \textbf{right adjoint} to \(L\), denoted \(L \dashv R\), if there is a bijection:
  \begin{equation} \label{eq:adjunction-hom-iso}
    \phi_{X,Y}: \Hom_{\category{D}}(L(X), Y) \isoarrow \Hom_{\category{C}}(X, R(Y))
  \end{equation}
  which is natural in both \(X \in \category{C}\) and \(Y \in \category{D}\).
\end{definition}

\begin{definition}[Unit and counit]\label{def:unit-counit}
  The adjunction \eqref{eq:adjunction-hom-iso} is uniquely determined by two \hyperref[def:natural-transformation]{natural transformations}:
  \begin{itemize}
    \item The \textbf{unit} \(\eta: \id_{\category{C}} \implies R \circ L\), where \(\eta_{X} = \phi_{X, L(X)}(\id_{L(X)})\).
    \item The \textbf{counit} \(\varepsilon: L \circ R \implies \id_{\category{D}}\), where \(\varepsilon_{Y} = \phi^{-1}_{R(Y), Y}(\id_{R(Y)})\).
  \end{itemize}
  These satisfy the \textbf{triangular identities}: \(R(\varepsilon_{Y}) \circ \eta_{R(Y)} = \id_{R(Y)}\) and \(\varepsilon_{L(X)} \circ L(\eta_{X}) = \id_{L(X)}\).
\end{definition}

\begin{remark}
  \hyperref[def:adjoint-functors]{Left and right adjoints} are \hyperref[def:dual]{dual concepts}.
\end{remark}


\begin{proposition}[Adjoints are unique]\label{prop:uniqueness-adjoints}
  The \hyperref[def:adjoint-functors]{right adjoint} (resp. \hyperref[def:adjoint-functors]{left adjoint}) to a \hyperref[def:functor]{functor}, provided it exists, is unique up to \hyperref[def:natural-transformation]{natural isomorphism}
\end{proposition}
\begin{proof}
  Let \(L: \category{C} \to \category{D}\) be a \hyperref[def:functor]{functor}.
  
  Let \(R_{1}, R_{2}: \category{D} \to \category{C}\) be two \hyperref[def:functor]{functors} which are \hyperref[def:adjoint-functors]{right adjoint} to \(L\).
  This means that for all \(X \in \category{C}\) and \(Y \in \category{D}\) we have bijections
  \begin{align*}
    \phi_{X, Y}:& \Hom_{\category{D}}(L(X), Y) \isoarrow \Hom_{\category{C}}(X, R_{1}(Y)) \\
    \psi_{X, Y}:& \Hom_{\category{D}}(L(X), Y) \isoarrow \Hom_{\category{C}}(X, R_{2}(Y))
  \end{align*}
  Using this we can define a bijection
  \begin{equation*}
    \Theta_{X, Y} = \psi_{X, Y} \circ \phi_{X, Y}^{-1}: \Hom_{\category{C}}(X, R_{1}(Y)) \isoarrow \Hom_{\category{C}}(X, R_{2}(Y))
  \end{equation*}
  Since \(\psi\) and \(\phi\) are natural in \(X\) so is \(\Theta\).
  Therefore for a fixed \(Y \in \category{D}\) we have a \hyperref[def:natural-transformation]{natural isomorphism}
  \begin{center}
    \begin{tikzcd}[column sep=huge]
      \category{C} \arrow[r, bend left=50, "{\Hom_{\category{C}}(-, R_{1}(Y))}" ""{name=source, below}]
      \arrow[r, bend right=50, "{\Hom_{\category{C}}(-, R_{2}(Y))}"' ""{name=target, above}]
      & \category{Set}
      \arrow[Rightarrow, from=source, to=target, "\Theta_{-, Y}"]
    \end{tikzcd}
  \end{center}

\end{proof}



\begin{proposition}\label{prop:adjoints-preserve-limits}
  \hyperref[def:adjoint-functors]{Right adjoints} preserve all \hyperref[def:categorical-limit]{limits} that exist in \(\category{D}\).
  
  \hyperref[def:dual]{Dually} \hyperref[def:adjoint-functors]{left adjoints} preserve all \hyperref[def:categorical-colimit]{colimits} that exist in \(\category{C}\).
\end{proposition}


\section{Functor categories}

\begin{definition}[Functor category]\label{def:functor-category}
  Let \(\category{C}\) and \(\category{D}\) be \hyperref[def:big-category]{categories}. The \textbf{functor category} \([C, D]\) (also denoted \(\category{D}^{\category{C}}\)) has:
  \begin{itemize}
    \item \textbf{Objects}: \hyperref[def:functor]{functors} \(F: \category{C} \to \category{D}\).
    \item \textbf{Morphisms}: \hyperref[def:natural-transformation]{natural transformations} \(\eta: F \implies G\).
  \end{itemize}
  Composition and identities are defined componentwise.
\end{definition}

\begin{lemma}
  Let \(\category{C}\) be a \hyperref[def:small-category]{small category}.
  Then
  \begin{itemize}
    \item If \(\category{D}\) is a \hyperref[def:category]{locally small category} then the \hyperref[def:functor-category]{functor category} \([\category{C}, \category{D}]\) is \hyperref[def:category]{locally small}.
    \item If \(\category{D}\) is a \hyperref[def:small-category]{small category} then the \hyperref[def:functor-category]{functor category} \([\category{C}, \category{D}]\) is \hyperref[def:small-category]{small}.
  \end{itemize}
\end{lemma}


\begin{lemma}
  Let \(\category{C}\) and \(\category{D}\) be two \hyperref[def:big-category]{categories} and \([\category{C}, \category{D}]\) be a \hyperref[def:functor-category]{functor category}.

  If \(\category{D}\) has \hyperref[def:categorical-limit]{limits} (resp. \hyperref[def:categorical-colimit]{colimits}) of shape \(\category{J}\) then \([\category{C}, \category{D}]\) has  \hyperref[def:categorical-limit]{limits} (resp. \hyperref[def:categorical-colimit]{colimits}) of shape \(\category{J}\).
  Furthermore, they are computed objectwise.
\end{lemma}


\begin{corollary}[{Yoneda embedding - \cite[Corollary 2.2.8]{riehlCategoryTheoryContext}}]
  Let \(\category{C}\) be a \hyperref[def:category]{locally small category}.
  The \hyperref[def:functor]{functor} \(Y: \opcategory{C} \to [\category{C}, \Set]\) given by \(X \mapsto \Hom(X, -)\) is \hyperref[def:functor-types]{fully faithful}. This allows us to view any locally small category as a \hyperref[def:subcategory]{full subcategory} of a \hyperref[def:functor-category]{functor category}.
\end{corollary}

\section{Higher categories}

We can go a step further.
\begin{definition}\label{def:natural-transformation-category}
  Let \(F: \category{C} \to \category{D}\), \(G: \category{C} \to \category{D}\) be two \hyperref[def:functor]{functors}.
  The \textbf{natural transformations category} \(\category{Nat}(F, G)\) constists of:
  \begin{itemize}
    \item \textbf{Objects}: \hyperref[def:natural-transformation]{natural transformations} \(F \Rightarrow G\).
    \item \textbf{Morphisms}: 
  \end{itemize}
\end{definition}



\end{document}
